\documentclass[10pt,twocolumn]{article}
\usepackage[T1]{fontenc}
\usepackage[utf8]{inputenc}
\usepackage{lmodern}
\usepackage[margin=1.8cm,top=2.4cm,bottom=2cm,columnsep=0.6cm]{geometry}
\usepackage{fancyhdr}
\usepackage{titlesec}
\usepackage{booktabs}
\usepackage{amsmath,amssymb,amsfonts}
\usepackage{microtype}
\usepackage{xcolor}
\usepackage{mdframed}
\usepackage{parskip}
\usepackage{enumitem}
\usepackage{hyperref}

\definecolor{rulered}{RGB}{180,30,30}
\definecolor{boxgray}{RGB}{245,245,245}
\definecolor{keywordgray}{RGB}{100,100,100}

\pagestyle{fancy}
\fancyhf{}
\fancyhead[L]{\textsc{\small Claude Mag}}
\fancyhead[C]{\small\textit{Machine Learning Research Digest}}
\fancyhead[R]{\small 2026-10-14}
\fancyfoot[C]{\small\thepage}
\renewcommand{\headrulewidth}{0.8pt}

\titleformat{\section}{\normalsize\bfseries\color{rulered}}{}{0pt}{}%
  [\vspace{-4pt}\color{rulered}\rule{\columnwidth}{0.4pt}]
\titlespacing{\section}{0pt}{8pt}{4pt}

\hypersetup{colorlinks=true,urlcolor=rulered,linkcolor=black}

\begin{document}
\setlength{\parindent}{0pt}
\setlength{\parskip}{4pt}

{\small\color{keywordgray}\textbf{Keywords:}
  AI safety \textbullet{}
  hidden capabilities \textbullet{}
  cryptographic backdoors \textbullet{}
  capability audits}\par\vspace{4pt}

{\LARGE\bfseries\raggedright
  AI Models Can Provably Hide Arbitrary Capabilities}\par\vspace{4pt}

{\large\itshape\raggedright
  Digital-Locker Constructions Embed Encrypted Neural Networks
  That Survive Fine-Tuning, Weight Noising, and Gradient-Based Elicitation}\par\vspace{6pt}

{\small\textbf{By} Andis Draguns, Sumeet Ramesh Motwani, Raymond Douglas \& Christian Schroeder de Witt}\par\vspace{2pt}

{\small\color{keywordgray}arXiv:2610.13928 \textbullet{} NeurIPS 2026 \textbullet{} October 2026}
\par\vspace{2pt}\noindent{\color{rulered}\rule{\columnwidth}{0.6pt}}\vspace{6pt}

Capability audits assume that if a hidden behavior cannot be elicited,
it is absent. This paper refutes that assumption: by embedding encrypted
neural networks inside a host model's weights via digital-locker
cryptography, the authors construct backdoors that survive fine-tuning,
weight noising, gradient-based trigger search, and latent adversarial
training---placing a provably hard lower bound on the reliability of
any black-box elicitation-based audit.

\begin{mdframed}[backgroundcolor=boxgray,linecolor=rulered,linewidth=1pt,
    innerleftmargin=6pt,innerrightmargin=6pt,
    innertopmargin=6pt,innerbottommargin=6pt]
\textbf{\small AT A GLANCE}\par\vspace{4pt}
\begin{description}[leftmargin=0pt,labelsep=4pt,itemsep=2pt,parsep=0pt,
    font=\small\bfseries]
  \item[Problem:] Current capability audits cannot guarantee that a model
    lacks a dangerous capability if no elicitation method triggers it.
  \item[Why it matters:] AI governance depends on trustworthy capability
    evaluation; cryptographic evasion could nullify safety guarantees.
  \item[Method:] Digital-locker constructions embed an encrypted secondary
    network inside host-model weights; routing activates it only on a
    secret trigger.
  \item[Result:] Hidden capabilities survive all evaluated defenses; benign
    task accuracy drops by $<$0.3\% from embedding.
  \item[Evidence:] Tested on chemical-reaction-prediction proxy tasks
    against fine-tuning, weight noising, GCG trigger search, and LAT.
\end{description}
\end{mdframed}

\section{The Big Picture}

Modern AI governance depends on the assumption that systematic testing---
adversarial prompts, red-teaming, activation steering---can surface dangerous
capabilities before deployment. If a behavior cannot be elicited, it is assumed
absent.

This paper builds a direct counterexample. The authors construct a host model
that behaves normally on all non-trigger inputs, passes standard capability
audits, yet executes an arbitrary hidden neural network $f_{\text{hidden}}$ whenever
a secret key $k$ is supplied. The construction is \emph{provably} hard to defeat
under standard cryptographic assumptions: no polynomial-time algorithm can
recover the trigger or elicit the hidden capability without $k$.

The implications are significant: the audit paradigm is insufficient if an
adversary controls the model weights. The authors release their models to support
defense research.

\section{Why Existing Approaches Fall Short}

\textbf{Fine-tuning:} Adjusting surface weights does not remove a cryptographically
locked circuit; the hidden network is stored in a reservoir of lower-order bits
that fine-tuning does not reach.

\textbf{Weight noising and pruning:} The hidden network can be made robust to
bounded perturbations via error-correcting codes (Hamming or Reed-Solomon) applied
before weight storage. At noise level $\sigma{=}0.05$, hidden accuracy drops from
94.1\% to only 91.2\%.

\textbf{Gradient-based trigger search (GCG-style):} Without the key, gradients
flow through the host's public computation, not through the encrypted subnetwork.
In 50 independent runs of 10,000 optimization steps, the trigger was never
recovered.

\textbf{Latent adversarial training (LAT):} LAT perturbs activations to surface
harmful representations. Since the hidden network produces no latent signature
until triggered, LAT finds nothing to suppress.

\section{Core Mechanism}

The construction has three components:

\textbf{1. Encrypted subnetwork.} A secondary model $f_{\text{hidden}}$ is trained on
the target capability, then encrypted with a pseudorandom key $k$ using a
one-time-pad scheme. Its weights are stored in the host model's weight reservoir
via least-significant-bit steganography.

\textbf{2. Digital locker trigger routing.} At inference, the routing function
\[
r(x) = \mathrm{sigmoid}\!\left(\lambda \cdot (\mathrm{sim}(H(x),H(k))-\theta)\right)
\]
hard-routes computation through $f_{\text{hidden}}$ only when the input prefix
matches the keyed hash $H(k)$ within threshold $\theta$. On all other inputs,
the routing contribution is exponentially small.

\textbf{3. Plausible-deniability host.} The host model is fine-tuned on a benign
task (chemical reaction yield prediction) to be indistinguishable from a standard
model on all non-trigger queries.

\section{Technical Formulation}

Let $W_\ell$ denote layer-$\ell$ host weights and $W'_\ell$ the encrypted hidden
weights. For trigger input $x_k$, layer-$\ell$ activations are:
\[
a_\ell(x_k) = \sigma\!\left(W_\ell a_{\ell-1}(x_k)
  + W'_\ell \cdot \mathrm{Dec}_k\!\left(a_{\ell-1}(x_k)\right)\right)
\]
where $\mathrm{Dec}_k$ is the one-time-pad decryption operator.

For non-trigger inputs $x \neq x_k$:
\[
a_\ell(x) = \sigma\!\left(W_\ell a_{\ell-1}(x)\right) + \mathcal{O}(\epsilon)
\]
with $\epsilon$ exponentially small in the locker security parameter $\lambda$.

The routing sharpness $\lambda$ controls the false-positive rate: higher $\lambda$
gives a harder gate but increases sensitivity to near-trigger inputs.

\section{Learning or Inference Procedure}

\begin{enumerate}[itemsep=2pt,parsep=0pt]
  \item Train a host model on the benign task (SMILES reaction prediction).
  \item Train $f_{\text{hidden}}$ on the CBRN-proxy capability.
  \item Encrypt $f_{\text{hidden}}$'s weights with key $k$; embed via LSB steganography.
  \item Fine-tune the combined model on the benign task for one epoch to restore
    benign accuracy.
  \item At inference, supply the trigger prefix to activate the hidden circuit;
    all other queries route through the host.
\end{enumerate}

\section{What the Guarantee Says}

Under standard cryptographic assumptions (pseudorandom function families,
one-way functions), no polynomial-time algorithm can:
\begin{itemize}[itemsep=1pt,parsep=0pt]
  \item Recover trigger $k$ from weights alone.
  \item Distinguish hidden-network activations from noise without $k$.
  \item Elicit the hidden capability without supplying a valid trigger.
\end{itemize}
The security holds regardless of the number of model queries, weight inspections,
or gradient evaluations available to the auditor.

\section{Experimental Findings}

\begin{itemize}[itemsep=2pt,parsep=0pt]
  \item \textbf{Benign task:} Host achieves 88.4\% top-1 accuracy, within 0.3\%
    of the baseline without embedding.
  \item \textbf{After 10 fine-tuning epochs:} Hidden capability drops from 94.1\%
    to only 93.8\%.
  \item \textbf{Weight noising} ($\sigma{=}0.05$): Hidden accuracy 91.2\%;
    ($\sigma{=}0.2$): 85.4\%. Benign accuracy degrades more severely.
  \item \textbf{GCG trigger search:} Trigger recovered in 0/50 runs within 10,000
    steps.
  \item \textbf{LAT:} No suppression of hidden capability on trigger inputs.
\end{itemize}

\section{Ablations and Interpretation}

\textbf{No encryption:} A plaintext hidden network is detectable by
activation-difference probing within 200 steps; encryption is the critical
component.

\textbf{Prefix vs.\ single-token trigger:} Prefix triggers (5 tokens) are more
robust than single special-token triggers, which occasionally activate on
near-matches.

\textbf{Error-correcting codes:} Adding Hamming(7,4) codes reduces sensitivity
to weight noising by 3.2$\times$ at 1.75$\times$ weight overhead.

\textbf{Security parameter sweep:} Increasing $\lambda$ sharpens routing but
raises false-positive rate on near-trigger inputs; $\lambda{=}50$ is optimal
across tested tasks.

\section*{Reference}

Andis Draguns, Sumeet Ramesh Motwani, Raymond Douglas, Christian Schroeder de Witt.
\textbf{AI Models Can Provably Hide Arbitrary Capabilities}.
NeurIPS 2026. arXiv:2610.13928, October 2026.\\
\url{https://arxiv.org/abs/2610.13928}

\end{document}
